\subsection{Baseline reproduced, then beaten (\eb{})}
\label{sec:baseline}

\textbf{The baseline.} NRMS \cite{wu2019nrms} is a neural recommender with two parts: a news
encoder that turns each article into a vector, and a user encoder that uses attention over the
user's clicked articles to decide which past clicks matter for the current prediction. We ran the
organisers' \texttt{ebnerd\_nrms\_docvec.py} unmodified on \eb{} small with the provided document
vectors, history size 20, 4 negatives per positive, batch 32 and seed 16. It reached val AUC
\textbf{0.6484} (best epoch 3, early-stopped at 7, 54 minutes on 16 CPU cores of a cluster node).

\textbf{That number is not comparable to ours.} The benchmark script trains on the organisers'
train and validation blocks together and validates on its own last day, which in our split means
it trains on our val and test impressions. So we also drove the same model, settings, vectors and
data loader from our temporal split: train on our train, stop early on our val, score our test
once. That run is the baseline we compare against.

\begin{table}[tb]
\centering\small
\begin{tabular}{@{}lcc@{}}
\toprule
\eb{} small, our temporal split & val AUC & test AUC \\
\midrule
NRMS, benchmark script and protocol (not comparable) & 0.6484 & -- \\
NRMS baseline, history size 20 & 0.5969 & 0.5938 \\
NRMS improved, history size 50 & 0.6042 & 0.6032 \\
\quad improved $-$ baseline & \dcell{$+$0.0072}{[$+$0.0058, $+$0.0088]} & \dcell{$+$0.0094}{[$+$0.0086, $+$0.0102]} \\
\midrule
our two-stage re-ranker & \EBFinalVal & \EBFinalTest \\
\bottomrule
\end{tabular}
\caption{NRMS reproduced and beaten with one change, paired bootstrap. Both NRMS arms score the
identical 64,365 val and 244,647 test impressions, aligned by impression id.}
\label{tab:nrms}
\end{table}

\textbf{The one principled change: a longer history, 20 to 50 clicks.} The motivation is our own
measurement. On clean data, the mean-of-history user vector improves steadily as the history
grows, from 0.5416 at 10 clicks to 0.5545 at 100 (Appendix~\ref{app:sweeps}). NRMS builds its user
representation from exactly that history, so it should benefit too. It does: the gain is
significant on both splits and larger on test than on val, the opposite of what over-fitting to
val would produce. We stopped at 50 because attention cost grows with the square of the history
length. A guard asserts every article id reaches the model (coverage 1.0000 on all four checks),
since a silent id mismatch would train on zero vectors and still report a believable AUC.

\textbf{Why our system beats NRMS by about 0.19 AUC.} Not because trees beat attention: NRMS reads
article content and click order only, and has no click popularity, our most valuable signal. Our
own content-only arm, \feat{stage1_only}, scores 0.5498 val, close to NRMS. The comparison is
between a model that sees the click log's behavioural signal and one that does not.
